\begin{savequote}
\begin{verbatim}
 _______________________________
( Don't look back, the lemmings )
( are gaining.                  )
 -------------------------------
       o   ,__,
        o  (oo)____
           (__)    )\
              ||--|| *

\end{verbatim}
\end{savequote}

\chapter
[\CO{[TAE]} Tätigkeitsbild und Orientierung]
{\CC{[TAE]}\\Tätigkeitsbild und Orientierung}
\label{chp:TAE}

\minitoc

\section{Präklinik}

\subsection{Sanitäter}

\subsubsection{Sanitäter in der Theorie}

Der Beruf bzw. die Tätigkeit des Sanitäters wurde 2002 durch das neue 
\E{Sanitätergesetz} (SanG) neu geregelt. Demnach unterteilt man in
\E{Rettungssanitäter} oder \E{Notfallsanitäter}. Der Beruf bzw. die Tätigkeiten des Sanitäters dürfen nur nach Maßgabe des Sanitätergesetzes
 ausgeübt werden. 

Sanitäter haben ihre Tätigkeit \E{ohne Ansehen der Person} gewissenhaft
auszuüben. Sie haben das Wohl der Patienten und der betreuten Personen nach
\E{Maßgabe der fachlichen und wissenschaftlichen Erkenntnisse und Erfahrungen}
zu wahren. 
% Nötigenfalls ist ein Notarzt oder, wenn ein solcher nicht zur Verfügung steht,
% ein sonstiger zur selbständigen Berufsausübung berechtigter Arzt anzufordern.

\Es{Fortbildung} ist ein wichtiger Teil der Tätigkeit, Sanitäter haben sich
laufend tätigkeitsrelevant fortzubilden. Darüber hinaus muss er \E{Berufs- bzw.
Tätigkeitspflichten} erfüllen (Dokumentationspflicht, Verschwiegenheitspflicht,
Auskunftspflicht).


\Layout{60}{Allgemeiner Tätigkeitsbereich}{%
Der Tätigkeitsbereich des Sanitäters beinhaltet eine Reihe von Tätigkeiten,
welche \Es{eigenverantwortlich} durchgeführt werden. Dazu gehört die
eigenverantwortliche Anwendung von Maßnahmen der qualifizierten Ersten Hilfe,
Sanitätshilfe und Rettungstechnik, einschließlich diagnostischer und
therapeutischer Verrichtungen.
}
{}
{./images/TAE/dif-w2-zelt-1.jpg}
{Fürsorge für erkrankte oder
hilfsbedürftige Personen: Die Kernaufgabe eines jeden Sanitäters.}
{}


\Layout{60}{Tätigkeitsbereich des Rettungssanitäters}{%
Der Tätigkeitsbereich des \Es{Rettungssanitäters} umfasst die selbständige und
\E{eigenverantwortliche Versorgung und Betreuung kranker, verletzter und
sonstiger hilfsbedürftiger Personen}, die medizinisch indizierter Betreuung
bedürfen, vor und während des Transports, die Übernahme sowie die Übergabe des
Patienten oder der betreuten Person im Zusammenhang mit einem Transport, 
eine qualifizierte Durchführung von lebensrettenden Sofortmaßnahmen sowie  
die sanitätsdienstliche Durchführung von Sondertransporten.
}
{}
{./images/TAE/tae-teamwork-transfer-lbk.jpg}
{Auch die gute Zusammenarbeit mit dem
Spitalspersonal ist wichtig, wie zum Beispiel bei dieser Überstellung eines
Intensivpatienten. Wenn alle Beteiligten danach so fröhlich sind ist das ein
gutes Zeichen.}
{}



\Layout{2}{Tätigkeitsbereich des Notfallsanitäters}{%
Der Tätigkeitsbereich des \Es{Notfallsanitäters} umfasst zusätzlich die
\E{Unterstützung des Arztes} bei allen notfall- und katastrophenmedizinischen
Maßnahmen einschließlich der Betreuung und des sanitätsdienstlichen Transports von
\E{Notfallpatienten}, 
die Verabreichung von dafür freigegebenen
\E{Arzneimitteln}, 
die eigenverantwortliche Betreuung der berufsspezifischen Geräte, Materialien und
Arzneimittel und die Mitarbeit in der \E{Forschung}. Notfallsanitäter können
die Berechtigung zur Durchführung allgemeiner und besonderer
\E{Notfallkompetenzen} erwerben (Arzneimittellehre, Venenzugang und Infusion,
Beatmung und Intubation). 
}{}{}{}{}


\Layout{60}{Dienstverhältnis}{%
Tätigkeiten des Sanitäters dürfen \E{ehrenamtlich}, \E{berufsmäßig}, als
\E{Zivildienstleistender} oder unter anderen bestimmten Umständen\ZFootnote{Die
Tätigkeiten des Sanitäters dürfen ferner als Soldat im Bundesheer, als Organ des
öffentlichen Sicherheitsdienstes, Zollorgan, Strafvollzugsbediensteter,
Angehöriger eines sonstigen Wachkörpers} ausgeübt werden. Die \Es{berufsmäßige}
Ausübung von Tätigkeiten des Sanitäters setzt neben der jeweiligen
Fachausbildung die Absolvierung eines \E{Berufsmoduls} voraus. 
}
{}
{./images/TAE/tae-okt26-alle.jpg}
{Im Notfall ist die
Organisationszugehörigkeit nebensächlich.}
{}


\Layout{2}{Befristung}{%
Die Berufs- und Tätigkeitsberechtigung ist mit jeweils zwei Jahren
\E{befristet}. Zur Verlängerung der Berufs- und Tätigkeitsberechtigung bedarf
es der Absolvierung von Fortbildungen sowie einer Rezertifizierung.
}{}{}{}{}


% (2) Notfallpatienten gemäß Abs. 1 Z 2 sind Patienten, bei denen im Rahmen einer
% akuten Erkrankung, einer Vergiftung oder eines Traumas eine lebensbedrohliche
% Störung einer vitalen Funktion eingetreten ist, einzutreten droht oder nicht
% sicher auszuschließen ist.

\subsubsection{Sanitäter in der Praxis}

Es ist derzeit noch relativ unklar wie der Sanitäter in das etablierte
Berufsfeld der Gesundheitsberufe "`hineinpasst"'. Die Berufsgruppen der Ärzte,
Hebammen, Pflegepersonal usw. gibt es schon seit Jahrzehnten oder gar
Jahrhunderten und es blieb genug Zeit für Streitigkeiten bezüglich der
jeweiligen Kompetenz und der Abgrenzung der Tätigkeitsfelder. Des öfteren hat
der Gesetzgeber Regelungen aufgestellt und die Gerichte Orientierungspunkte
gegeben; darüber hinaus wurden durch die ständige Arbeit miteinander Realitäten
"`geschaffen"'. In manchen Punkten aber auch heutzutage einiges Unklar oder
einem stetigen Wandel unterworfen. 

\Layout{2}{Sanitätergesetz 2002}{%
Mit der Schaffung des Sanitätergesetzes im Jahr 2002 wollte man die Tätigkeit
der Sanitäter endlich klar regeln und Rechtssicherheit schaffen. Gleich vorweg,
wirkliche Klarheit über die Tätigkeiten des Sanitäters wurde nicht geschaffen,
wie unzählige Diskussionen zu diversen Themen beweisen. Bezüglich der genauen
rechtlichen Regelungen sei auf das Kapitel \prettyref{chp:JUS} verwiesen, an
dieser Stelle wollen wir uns aus praktischer Sicht dem Tätigkeitsfeld und den
Aufgaben der Sanitäter widmen. 
}{}{}{}{}

\Layout{2}{Das Miteinander}{%
Es scheint ein Naturgesetz zu sein dass, sobald zwei oder mehr Berufsgruppen
aufeinander treffen, es zu Streitigkeiten bezüglich der jeweiligen Kompetenzen,
Rechte und Befugnisse kommt. Auf der Strecke bleibt oft das Potential welches
entwickelt werden könnte, würden die Berufsgruppen konstruktiv und
zielgerichtet zusammenarbeiten. In diesem Sinne sind die \AASS{} auf ein
interdisziplinäres Miteinander der Berufsgruppen.}{}{}{}{}

\opt{STATUS-DRAFT}{
\Layout{2}{Der Sanitäter -- früher}{%
\TakeNotesLarge
}{}{}{}{}
}

\Layout{2}{Der Sanitäter -- heute}{%
Das Aufgabengebiet des Sanitäter von heute wird immer mehr um Tätigkeiten
erweitert, die medizinisch-wissenschaftliche Kenntnisse voraussetzen
und noch vor einiger Zeit nur den Ärzten oder anderen Berufsgruppen
vorbehalten waren. 

Dies geschah anfangs durch das SanG, später aber auch durch die Anforderungen,
welche sich zwangsläufig aus der Arbeit in der Praxis ergeben und z\,T. gar
nicht gesetzlich geregelt sind. 

In letzter Zeit kommt es auch zu einem Einfluss privater, meist kommerzieller,
internationaler Ausbildungsprogramme wie z.\,B.
\Abk{AMLS}{\texttrademark}\footnote{Advanced Medical Life
Support{\texttrademark}}, \Abk{ITLS}{\texttrademark}\ZFootnote{International
Trauma Life Support{\texttrademark}}, \Abk{PHTLS}{\texttrademark}\ZFootnote{Prehospital Trauma Life Support{\texttrademark}} etc. Deren Lehrinhalte sind meist auf den Kompetenzbereich von
Paramedics\ZFootnote{Paramedics sind gutausgebildete Sanitäter im
englischsprachigen Raum, welche in nicht-Notarzt-basierten Systemen eingesetzt
werden, d.\,h. sie haben in diesen Ländern das Personal mit der höchsten
medizinischen Ausbildung, welches im präklinischen Regelbetrieb eingesetzt
wird.} und reizen Grenzen dessen, was allgemein als Aufgaben der Sanitäter
angesehen wird, weitgehend aus, oder überschreiten diese. Diesbezüglich kommt
es immer wieder zu Diskussionen, was Sanitäter tun bzw.
unterlassen\footnote{Schädliche Dinge zu unterlassen ist genauso wichtig wie
nützliche Dinge zu tun.} \sout{dürfen} \E{müssen}\footnote{Meist kann man es
sich nicht aussuchen, ob man eine Maßnahme ergreift oder nicht. Handlungen,
die man machen oder unterlassen darf und die das Patientenwohl fördern,
\E{müssen} i.\,d.\,R. entsprechend vorgenommen oder unterlassen werden.}, und was sie
\E{nicht} tun oder unterlassen dürfen. }{}{}{}{}

% \clearpage


\subsubsection{Anforderungen -- Beanspruchung -- Belastung: Was erwartet die
Welt von mir?}

\label{topic:belastung}

\Layout{2}{{\TxQuerverweise}}{%
\begin{inparaitem}
  \item Gesprächsführung: \CO{[AUP]} (\ref{chp:AUP});
  \item Belastungsreaktion, Burn-out:
  \prettyref{topic:ptss}, \prettyref{topic:burn-out}
\end{inparaitem}
}{}{}{}{}

\begin{description}
  \item[Anforderungen:] \H{fachliches Wissen --
  praktische Kompetenzen -- dieses in schwierigen Situationen umsetzen zu
  können professionelle Haltung -- social/soft skills zwischenmenschliche
  Fähigkeiten, Frustrationstoleranz -- Umgang mit mehrdeutigen
  Situationen bzw. "`schwierigem Klientel"' -- Teamdifferenzen -- \ldots}
\end{description}


\begin{center}
\ttfamily\small
\begin{minipage}{45em}
\begin{verbatim}
 _________________________________________
( Your object is to save the world, while )
( still leading a pleasant life.          )
 -----------------------------------------
  o
   o   \_\_    _/_/
    o      \__/
           (oo)\_______
           (__)\       )\/\
               ||----w |
               ||     ||

\end{verbatim}
\end{minipage}
\end{center}

\Layout{2}{Theoretische Anforderungen}{%
Die Tätigkeit im Rettungs- und Krankentransportwesen geht aufgrund der
Komplexität der Aufgaben mit einer Fülle von körperlichen, wissensmäßigen und
psychischen Anforderungen einher. Einen nur kargen Hinweis darüber, welche Qualitäten von
einem (Zivildienstleistenden) Sanitäter erwartet werden, liefert das
Anforderungsprofil des AMS, nämlich:



\begin{itemize}
\item Körperliche und psychische Belastbarkeit
\item rasches Auffassungs- und Reaktionsvermögen
\item Einfühlungsvermögen
\item Beobachtungsgabe
\item Kommunikations- und Teamfähigkeit
\item Zuverlässigkeit
\item Verantwortungsbewusstsein
\end{itemize}
}{%

}{}{}{}



Im Folgenden soll der Versuch unternommen
werden diese nüchternen Auflistung von Begrifflichkeiten
plastischer darzustellen und zu erweitern.

\Layout{0}{Anforderungen und Belastungen in der Praxis}{%
Es versteht sich von selbst, dass die professionelle Versorgung erkrankter bzw.
verunfallter Menschen ein medizinisches Wissen voraussetzt und damit auch
einhergehend die praktischen Kompetenzen, dieses \Es{zuverlässig und
eigenverantwortlich} in schwierigen Situationen umzusetzen zu können. 

Dringliche
Versorgungssituationen bzw. Notfälle, in denen jede zeitliche Verzögerung
zulasten des Patienten fallen könnten, machen ein \Es{schnelles Einschätzen der
Lage} so wie die Bereitschaft \Es{Verantwortung zu übernehmen} notwendig. Da es
sich stets um eine Interaktion zwischen zwei Personen -- nämlich Versorgendem
und zu Versorgenden -- handelt, gelten zwischenmenschliche Fähigkeiten wie
höfliche Umgangsformen, \Es{Kommunikationsfähigkeit} sowie \Es{Einfühlungsvermögen} als
Grundvoraussetzung. 

Erkrankung bzw. Unfall stellen für die Betroffenen in
häufigen Fällen \E{kein alltägliches Geschehen} dar, sondern haben oft einen
kritisch-krisenhaften, jedenfalls aber frustrierend-belastenden Charakter.
Deshalb soll an dieser Stelle die besondere Wichtigkeit einer \Es{einfühlenden
und verständnisvollen Grundhaltung} betont werden. Dies kann für den
professionellen Helfer auch damit einhergehen, \E{seine eigenen Bedürfnisse}
bzw. Empfindungen für den Moment \E{hintanzustellen}, also eine hohe
\Es{Frustrationstoleranz}
erforderlich machen. 

\Fazit{Der (richtige) Umgang mit Frustration ist wesentlich, um im Rettungs-
Krankentransportdienst auf Dauer "`zu überleben"'.}

Im Rettungs- und Krankentransportdienst beschäftigte Personen arbeiten
 die meiste Zeit in Teams, sodass Qualitäten wie Toleranz,
Kritikfähigkeit und Selbstdisziplin ("`\Es{Teamfähigkeit}"') ebenso von hoher
Bedeutung sind (Die Praxis zeigt, dass hier häufig viel Spielraum für
Verbesserungen besteht \Code{;-)}).

Weiters sei auf die Wichtigkeit der \Es{\E{körperlichen} Belastungsfähigkeit}
als erforderliches Kriterium hingewiesen: das Tätigkeitsfeld des Rettungssanitäters
umfasst diverse Aufgaben, die mit körperlichem Einsatz und dementsprechender
Anstrengung verbunden sind.

Zu guter Letzt und im besonderen sei auf die Bedeutung \Es{psychischer
Belastbarkeit} hingewiesen, \AuthNotes{die auch in den folgenden Kapiteln zum
Thema gemacht werden soll,} v.\,a. in Hinblick auf den Umgang mit psychischem
Stress: 

Sanitätsdienst zu leisten, bedeutet tägliche \E{Konfrontation} mit dem
persönlichen Leid von Mitmenschen, des weiteren die \E{Verantwortlichkeit} –
ggfs. auch noch unter Zeitdruck – darauf professionell zu reagieren und zusätzlich die
besonderen und generellen \E{Stressauslöser des Arbeitsbereichs} eines
Sanitäters zu meistern. Dies benötigt eine hinreichend funktionierende
psychische Verarbeitungsfähigkeit, um ein gesundes seelisches Gleichgewicht bzw. sein
persönliches Wohlbefinden wahren und pathologischen Entwicklungen (Burnout,
Belastungsreaktion, etc.) präventiv entgegenwirken zu können.
}{%
}{}{}{}
\begin{multicols}{2}

% \Layout{27}{}{
% 
% }
% {}
% {}
% {}
% {}
\begin{figure}[H]
\includegraphics[width=\linewidth]{./images/TAE/dif-nacht-atmosphaere.jpg}
\Captionof{figure}{\AuthNotesPicLegalOk{}{}Hitze, Dunkelheit und
kontrolliertes Chaos: Ambulanzdienst auf dem Donauinselfest.}
\end{figure}

% \Layout{27}{}{
% 
% }
% {}
% {}
% {}
% {}
\begin{figure}[H]
\includegraphics[width=\linewidth]{./images/TAE/erste-hilfe-dif-schild-stange-2.jpg}
\Captionof{figure}{\AuthNotesPicLegalOk{}{}Licht, Menschenmassen und noch mehr
Hitze: Auch das ist das Donauinselfest.}
\end{figure}




% \Layout{27}{}{
% 
% }
% {}
% {}
% {}
% {}
% \begin{figure}[H]
% \includegraphics[width=\linewidth]{./images/TAE/jacken-alt-grau.jpg}
% \Captionof{figure}{\AuthNotesPicLegalOk{}{}%
% % In den letzten zehn Jahren hat sich viel getan: 
% Die alten, grauen Jacken erscheinen heute wie ein Relikt aus grauer
% Vorzeit.}
% \end{figure}



% \Layout{27}{}{
% 
% }
% {}
% {}
% {}
% {}
\begin{figure}[H]
\includegraphics[width=\linewidth]{./images/TAE/tae-okt26-bundesheer1.jpg}
\Captionof{figure}{\AuthNotesPicLegalOk{}{}Integraler Bestandteil der zivilen
und militärischen Landesverteidigung: Der Bundesheersanitäter.}
\end{figure}


\begin{figure}[H]
\includegraphics[width=\linewidth]{./images/gabriel-sebastian-aass/IMG_1853_1.jpg}
\Captionof{figure}{Die Zusammenarbeit mit der Exekutive -- auch an
sozialen Brennpunkten -- ist Alltag.}
\end{figure}

% \Layout{27}{}{
% 
% }
% {}
% {}
% {}
% {}
\begin{figure}[H]
\includegraphics[width=\linewidth]{./images/TAE/tae-teamwork-arzt-sani.jpg}
\Captionof{figure}{\AuthNotesPicLegalOk{}{}Arzt und Sanitäter bilden eine
Einheit.}
\end{figure}
\end{multicols}